% GENERATED FILE. Edit canonical lessons, then run npm run generate:books.
% canonical-source-hash: fnv1a64:c2e972f6172d2356
% canonical-lessons: JA-C440-shinchou, JA-C440-ukkari, JA-C440-darashinai, JA-C440-kichinto, JA-C440-tekitou

\chapter{Careful and Careless}
\label{ch:ja-careful-and-careless}
\hlchaptermodality{440}

\begin{chapteropening}
\textbf{By the end of this chapter:} \emph{I can say whether someone is careful or careless.}

Complete the last lesson of chapter 440: I can say whether someone is careful or careless.
\end{chapteropening}

\section[\texorpdfstring{\ja{しんちょう} (shinchō)}{しんちょう (shinchō)}]{\ja{しんちょう} (shinchō) — cautious; prudent}

\label{lesson:JA-C440-shinchou}



\begin{quote}
Before the new one: say the Japanese for rough, then the Japanese for courageous.
\end{quote}

\subsection*{You'll want to know: \ja{しんちょう}}
\textbf{\ja{しんちょう}} — \emph{shinchō} — \textquotedblleft{}cautious; prudent\textquotedblright{}.

\begin{quote}
\textbf{\ja{ははは} \ja{うんてんが} \ja{しんちょうです。}} — \emph{haha wa unten ga shinchō desu} — \textquotedblleft{}My mother is a careful driver.\textquotedblright{}
\end{quote}

A na-adjective: \ja{な} before a noun.

\begin{cousinweb}
\textbf{\ja{しんちょう}} (\emph{shinchō}) is Sino-Japanese (kango): two Chinese characters read with their on'yomi, readings borrowed from Chinese. They are \emph{shin}, \textquotedblleft{}be careful\textquotedblright{} (Middle Chinese \emph{dzyinH}) and \emph{chō}, \textquotedblleft{}important\textquotedblright{} (\emph{drjowngX}) (KANJIDIC2; Middle Chinese in Baxter's transcription).
\end{cousinweb}

\begin{grammarlens}[title={What you've built}]
One more word for describing what you lived and felt.
\end{grammarlens}

\subsection*{Your turn}
\begin{itemize}
\raggedright
  \item \emph{Say it:} \emph{shinchō}
  \item \emph{Say it:} \emph{shinchō}, once more
  \item \emph{Say it:} \emph{haha wa unten ga shinchō desu} — My mother is a careful driver.
  \item \emph{Recall:} say the Japanese for rough, then the Japanese for courageous, then say \emph{shinchō} again
\end{itemize}

\begin{tcolorbox}[breakable,colback=teal!4,colframe=teal!35!black,title={Before you move on}]
What does \ja{しんちょう} mean? (\textquotedblleft{}Cautious; prudent\textquotedblright{}.) Say it once more.
\end{tcolorbox}
\section[\texorpdfstring{\ja{うっかり} (ukkari)}{うっかり (ukkari)}]{\ja{うっかり} (ukkari) — thoughtlessly; inadvertently}

\label{lesson:JA-C440-ukkari}



\begin{quote}
Before the new one: say the Japanese for courageous, then the Japanese for cautious.
\end{quote}

\subsection*{You'll want to know: \ja{うっかり}}
\textbf{\ja{うっかり}} — \emph{ukkari} — \textquotedblleft{}thoughtlessly; inadvertently\textquotedblright{}.

\begin{quote}
\textbf{\ja{うっかり} \ja{かさを} \ja{わすれました。}} — \emph{ukkari kasa o wasuremashita} — \textquotedblleft{}Without thinking, I left my umbrella behind.\textquotedblright{}
\end{quote}

\begin{cousinweb}
\textbf{\ja{うっかり}} (\emph{ukkari}) is a native Japanese mimetic word, inherited rather than borrowed from Chinese, and its older origin is not settled (JMdict).
\end{cousinweb}

\begin{grammarlens}[title={What you've built}]
One more word for describing what you lived and felt.
\end{grammarlens}

\subsection*{Your turn}
\begin{itemize}
\raggedright
  \item \emph{Say it:} \emph{ukkari}
  \item \emph{Say it:} \emph{ukkari}, once more
  \item \emph{Say it:} \emph{ukkari kasa o wasuremashita} — Without thinking, I left my umbrella behind.
  \item \emph{Recall:} say the Japanese for courageous, then the Japanese for cautious, then say \emph{ukkari} again
\end{itemize}

\begin{tcolorbox}[breakable,colback=teal!4,colframe=teal!35!black,title={Before you move on}]
What does \ja{うっかり} mean? (\textquotedblleft{}Thoughtlessly; inadvertently\textquotedblright{}.) Say it once more.
\end{tcolorbox}
\section[\texorpdfstring{\ja{だらしない} (darashinai)}{だらしない (darashinai)}]{\ja{だらしない} (darashinai) — sloppy; slovenly}

\label{lesson:JA-C440-darashinai}



\begin{quote}
Before the new one: say the Japanese for cautious, then the Japanese for thoughtlessly.
\end{quote}

\subsection*{You'll want to know: \ja{だらしない}}
\textbf{\ja{だらしない}} — \emph{darashinai} — \textquotedblleft{}sloppy; slovenly\textquotedblright{}.

\begin{quote}
\textbf{\ja{あには} \ja{ふくが} \ja{だらしないです。}} — \emph{ani wa fuku ga darashinai desu} — \textquotedblleft{}My elder brother dresses sloppily.\textquotedblright{}
\end{quote}

An i-adjective.

\begin{cousinweb}
\textbf{\ja{だらしない}} (\emph{darashinai}) is a native Japanese word (wago), inherited rather than borrowed from Chinese; JMdict writes it in kana (JMdict).
\end{cousinweb}

\begin{grammarlens}[title={What you've built}]
One more word for describing what you lived and felt.
\end{grammarlens}

\subsection*{Your turn}
\begin{itemize}
\raggedright
  \item \emph{Say it:} \emph{darashinai}
  \item \emph{Say it:} \emph{darashinai}, once more
  \item \emph{Say it:} \emph{ani wa fuku ga darashinai desu} — My elder brother dresses sloppily.
  \item \emph{Recall:} say the Japanese for cautious, then the Japanese for thoughtlessly, then say \emph{darashinai} again
\end{itemize}

\begin{tcolorbox}[breakable,colback=teal!4,colframe=teal!35!black,title={Before you move on}]
What does \ja{だらしない} mean? (\textquotedblleft{}Sloppy; slovenly\textquotedblright{}.) Say it once more.
\end{tcolorbox}
\section[\texorpdfstring{\ja{きちんと} (kichinto)}{きちんと (kichinto)}]{\ja{きちんと} (kichinto) — neatly; properly}

\label{lesson:JA-C440-kichinto}



\begin{quote}
Before the new one: say the Japanese for thoughtlessly, then the Japanese for sloppy.
\end{quote}

\subsection*{You'll want to know: \ja{きちんと}}
\textbf{\ja{きちんと}} — \emph{kichinto} — \textquotedblleft{}neatly; properly\textquotedblright{}.

\begin{quote}
\textbf{\ja{ふくを} \ja{きちんと} \ja{たたんで} \ja{ください。}} — \emph{fuku o kichinto tatande kudasai} — \textquotedblleft{}Please fold the clothes neatly.\textquotedblright{}
\end{quote}

\begin{cousinweb}
\textbf{\ja{きちんと}} (\emph{kichinto}) is a native Japanese mimetic word, inherited rather than borrowed from Chinese, and its older origin is not settled (JMdict).
\end{cousinweb}

\begin{grammarlens}[title={What you've built}]
One more word for describing what you lived and felt.
\end{grammarlens}

\subsection*{Your turn}
\begin{itemize}
\raggedright
  \item \emph{Say it:} \emph{kichinto}
  \item \emph{Say it:} \emph{kichinto}, once more
  \item \emph{Say it:} \emph{fuku o kichinto tatande kudasai} — Please fold the clothes neatly.
  \item \emph{Recall:} say the Japanese for thoughtlessly, then the Japanese for sloppy, then say \emph{kichinto} again
\end{itemize}

\begin{tcolorbox}[breakable,colback=teal!4,colframe=teal!35!black,title={Before you move on}]
What does \ja{きちんと} mean? (\textquotedblleft{}Neatly; properly\textquotedblright{}.) Say it once more.
\end{tcolorbox}
\section[\texorpdfstring{\ja{てきとう} (tekitō)}{てきとう (tekitō)}]{\ja{てきとう} (tekitō) — suitable; offhand}

\label{lesson:JA-C440-tekitou}



\begin{quote}
Before the new one: say the Japanese for sloppy, then the Japanese for neatly.
\end{quote}

\subsection*{You'll want to know: \ja{てきとう}}
\textbf{\ja{てきとう}} — \emph{tekitō} — \textquotedblleft{}suitable; offhand\textquotedblright{}.

\begin{quote}
\textbf{\ja{この} \ja{ばしょは} \ja{てきとうです。}} — \emph{kono basho wa tekitō desu} — \textquotedblleft{}This place is suitable.\textquotedblright{}
\end{quote}

A na-adjective: \ja{な} before a noun.

\begin{cousinweb}
\textbf{\ja{てきとう}} (\emph{tekitō}) is Sino-Japanese (kango): two Chinese characters read with their on'yomi, readings borrowed from Chinese. They are \emph{teki}, \textquotedblleft{}suitable\textquotedblright{} (Middle Chinese \emph{syek}) and \emph{tō}, \textquotedblleft{}appropriate\textquotedblright{} (\emph{tangH}) (KANJIDIC2; Middle Chinese in Baxter's transcription).
\end{cousinweb}

\begin{grammarlens}[title={What you've built}]
One more word for describing what you lived and felt.
\end{grammarlens}

\subsection*{Your turn}
\begin{itemize}
\raggedright
  \item \emph{Say it:} \emph{tekitō}
  \item \emph{Say it:} \emph{tekitō}, once more
  \item \emph{Say it:} \emph{kono basho wa tekitō desu} — This place is suitable.
  \item \emph{Recall:} say the Japanese for sloppy, then the Japanese for neatly, then say \emph{tekitō} again
\end{itemize}

\begin{tcolorbox}[breakable,colback=teal!4,colframe=teal!35!black,title={Before you move on}]
What does \ja{てきとう} mean? (\textquotedblleft{}Suitable; offhand\textquotedblright{}.) Say it once more.
\end{tcolorbox}
